\documentclass[12pt,a4paper]{article}
\usepackage{graphicx}
\usepackage{amsmath}
\usepackage{hyperref}
\usepackage{geometry}
\geometry{margin=2.5cm}
\emergencystretch=3em

\title{\textbf{CDFD Runtime Paper 09: Tri-Regime Bioenergetics and Life-Number Diagnostics}}
\author{Steve Bico Mujjabi, MD\\
\small Independent Researcher\\
\small Founder, Vura Labs\\
\small Kampala, Uganda\\
\small ORCID: \href{https://orcid.org/0009-0001-0556-5516}{0009-0001-0556-5516}}
\date{May 2026}

\begin{document}
\maketitle

\begin{abstract}
\noindent \textbf{Executive synthesis.} This paper connects CDFD Runtime biology
to the tri-regime Life Number $\Lambda$. The model tracks energy input,
electron transport, proton coupling, stabilization load, relaxation time, and
maintenance cost. The local OOL phase experiment supplies a candidate
simulation boundary near $\Lambda=1$. The correct claim is not that the runtime
has solved the origin of life. The correct claim is that it provides a
reproducible way to test life-like throughput thresholds under explicit model
assumptions.
\end{abstract}

\paragraph{Greek-symbol convention.}
Unless a manuscript states a tighter equation-local meaning, Greek symbols are local CDFD variables or coefficients, not universal constants. Here $\Phi$ denotes driving flux, $\Psi_s$ denotes the adaptive operating ratio, $\Omega$ denotes a domain or state space, and $\Lambda$ denotes the Life Number where used. The coefficients $\alpha$, $\beta$, and $\gamma$ denote accumulation or reinforcement, relaxation or decay, and diffusion, coupling, or redistribution. The symbols $\delta$ and $\epsilon$ denote perturbation or tolerance scales, while $\eta$, $\lambda$, $\mu$, $\nu$, $\rho$, $\sigma$, $\tau$, and $\phi$ denote local efficiency, rate, baseline, frequency, density or correlation, noise or variance, time-scale, and phase or field terms. Subscripts restrict any coefficient to the named pathway, tissue, domain, or experiment.


\begin{figure}[ht]
\centering
\fbox{\begin{minipage}{0.92\textwidth}
\centering
\textbf{Runtime evidence path}\\[0.4em]
Prior CDFD mathematics $\longrightarrow$ CDFL/runtime state $\longrightarrow$ bounded output $\longrightarrow$ falsification target.
\end{minipage}}
\caption{Conceptual schematic for the runtime papers. The engine translates the mathematical frame from the prior CDFD papers into executable CDFL/runtime diagnostics; candidate outputs remain tied to explicit tests before any stronger claim is made.}
\end{figure}

\section{Prior CDFD Basis}
This manuscript does not rederive the mathematical core of CDFD. It uses the
Mujjabi Adaptive Operating Ratio and the constrained-vacuum notation introduced
in Part I \cite{MujjabiI2026}, the torus-knot hierarchy and boundary-mode work
from the Part I sequence \cite{MujjabiVI2026,MujjabiIX2026}, the public balance
law developed in the Part I capstone \cite{MujjabiX2026}, the Part I transport
tests \cite{MujjabiXII2026}, and the Life Number and tri-regime biology bridge
from Part II \cite{MujjabiPartII2026}. The runtime papers therefore act as an
execution layer: they keep the mathematics connected to commands, outputs,
falsification tests, and domain-specific limits.


\section{Biology as Bounded Throughput}
In this runtime, life-like behavior is modeled as persistent organized flow.
The basic local operating ratio is still
\begin{equation}
    \Psi_s=\frac{\Phi}{C}S M_s.
\end{equation}
For origins-of-life questions, however, a system-wide viability number is also
necessary because local $\Psi_s$ can be high in one region and low elsewhere.

\section{The Life Number}
The Life Number is
\begin{equation}
    \Lambda =
    \frac{\Phi \sigma_e \sigma_p \tau_{\mathrm{relax}}}
    {S_{\mathrm{stab}}E_{\mathrm{maintenance}}},
\end{equation}
where $\Phi$ is energy input, $\sigma_e$ is electron-transport effectiveness,
$\sigma_p$ is proton-coupling or coherence effectiveness,
$\tau_{\mathrm{relax}}$ is relaxation time, $S_{\mathrm{stab}}$ is stabilization
or dissipation burden, and $E_{\mathrm{maintenance}}$ is maintenance energy.

The working labels are:
\begin{itemize}
    \item $\Lambda<1$: decay-dominated;
    \item $\Lambda\approx1$: near-critical proto-biological window;
    \item $\Lambda>1$: sustained throughput in the model.
\end{itemize}

\section{Runtime Implementation}
The implementation lives in \texttt{engine/origins\_of\_life.py} and is called
from \texttt{engine/biology.py}. The runtime stores channel summaries in
\texttt{state.meta} keys such as \texttt{ool\_C\_input},
\texttt{ool\_sigma\_e}, \texttt{ool\_sigma\_p},
\texttt{ool\_S\_stability}, \texttt{life\_number}, and
\texttt{throughput\_J}.

\section{Experiment: OOL Phase Boundary}
The OOL phase script sweeps a 50 by 50 parameter grid over energy input and
geochemical constraint. The public report describes candidate points near
$\Lambda=1$, and the selected HDF5 output stores the grid.

The useful result is a model diagnostic: the runtime can locate and save a
phase-boundary surface in the declared toy parameter space. The result still needs testing against laboratory chemistry, geochemical constraints, and independent
origin-of-life models before any empirical claim is made.

\section{Biology Discovery Scripts}
The broader biology experiments include:
\begin{itemize}
    \item \texttt{experiments/run\_biology\_discovery.py}: metabolic saturation,
    chronic drift, and tri-regime throughput;
    \item \texttt{experiments/run\_bio\_mujjabi\_tests.py}: parasitic-drain and
    cycling-retention tests;
    \item \texttt{experiments/run\_drug\_synergy\_predictions.py}: same-axis
    versus cross-axis perturbation toy trajectories.
\end{itemize}
The public \texttt{runtime/diagnostics.py} module exposes guarded aromatic
source-mix rows (Part II Paper 7), photochemical endpoint status (Paper 11),
and the Life Number supply guardrail. The primary export path is
\texttt{python cdfd.py diagnostics}; the optional Streamlit studio and
\texttt{cdfd.py demo origins\_of\_life --source-scenario <name>} surface the
same values.

These are model diagnostics and hypothesis generators. They are not care
instructions or validated use rules.

\section{Falsification Conditions}
The tri-regime model fails or requires revision if:
\begin{itemize}
    \item measured prebiotic systems show no reproducible threshold near the
    proposed $\Lambda$ proxy;
    \item $\sigma_e$ or $\sigma_p$ cannot be measured or approximated in a
    domain-specific way;
    \item the boundary is indistinguishable from a simpler energy-only model;
    \item outputs become non-finite under reasonable parameter stress.
\end{itemize}

\section{Conclusion}
Tri-regime bioenergetics gives CDFD a bridge from local operating ratios to
system-level viability. Its strength is not rhetorical certainty; its strength
is that it produces explicit variables, runnable scripts, saved outputs, and
failure conditions.
\begin{thebibliography}{9}
\bibitem{MujjabiI2026}
Steve Bico Mujjabi, MD. \emph{Vortex Stability in a Constrained Transport Vacuum and the Origin of the Fine-Structure Constant}. CDFD Part I Paper I, preprint, 2026.

\bibitem{MujjabiVI2026}
Steve Bico Mujjabi, MD. \emph{The Universal Torus Knot Hierarchy: $Z_n$ Symmetry, the Cinquefoil Family, and the General Structure of CDFT Particle Families}. CDFD Part I Paper VI, preprint, 2026.

\bibitem{MujjabiIX2026}
Steve Bico Mujjabi, MD. \emph{Even-$n$ Torus Knots in CDFT: The Exclusion of $T(2,2)$, the Extension to $T(2,2m)$, and the Anti-Phase Mass Scaling Law}. CDFD Part I Paper IX, preprint, 2026.

\bibitem{MujjabiX2026}
Steve Bico Mujjabi, MD. \emph{The Vacuum Equation of State: Deriving Torus Knot Self-Consistency from the Public CDFD Balance Law $\Psi = \Phi/C$}. CDFD Part I Paper X, preprint, 2026.

\bibitem{MujjabiXII2026}
Steve Bico Mujjabi, MD. \emph{Friction, Superconductivity, Plasma States, and Transport Mysteries: Observable Tests of Adaptive Constraint}. CDFD Part I Paper XII, preprint, 2026.

\bibitem{MujjabiPartII2026}
Steve Bico Mujjabi, MD. \emph{CDFD Part II: Origins of Life and Tri-Regime Bioenergetics}. Public release archive, 2026, \url{https://doi.org/10.5281/zenodo.20264779}.
\end{thebibliography}
\end{document}
